\subsection{Generating the JAX program}
\label{sec:pfit-jax}

After the session passes the setup checks, the user runs

\begin{verbatim}
/pfit-jax <session-name>
\end{verbatim}

This step converts the checked model into the executable JAX program used for
fitting. The input is the readable session model,
\texttt{generated/user\_model.py}, together with the checked configuration in
\texttt{inputs/user\_input.yaml}. The output is
\texttt{generated/generated\_script.py}, which is the file imported by the
optimizer.

The translation is mechanical. It preserves the equations, observables, loss,
parameter bounds, initial conditions, dataset mappings, and solver choices that
were established during setup and checking. It does not introduce a new model,
reinterpret data columns, or change the fitted objective. Its role is to express
the validated problem in JAX-compatible form so that the objective can be
compiled, differentiated, and evaluated repeatedly during optimization.

For the two-experiment example, the translated program evaluates the same ODE
system for each dataset, applies the appropriate initial condition and forcing
history, computes the declared observable, and evaluates the masked residuals
specified in the loss. Blank measurements remain excluded from the objective.
If an integration fails for a trial parameter vector, the generated program
returns the configured failure loss rather than an invalid numerical value.

After writing the JAX program, the agent imports it in the session environment.
This import check catches syntax errors, missing symbols, unsupported numerical
operations, and invalid solver settings before the fitting run begins. The
script is then stamped with hashes of the configuration and source model from
which it was generated. Later run checks use this stamp to detect whether the
YAML or model file changed after translation. If either file is edited, the user
should rerun \texttt{/pfit-check} and \texttt{/pfit-jax} before fitting.

On success, the session contains a checked specification and an executable JAX
program. The user can then start parameter estimation with

\begin{verbatim}
/pfit-run <session-name>
\end{verbatim}
